% ---------------------------------------------------------------------------
% Capítulo 1 --- Descripción de la empresa
% ---------------------------------------------------------------------------
% La norma 3.13 exige que el cuerpo del trabajo incluya una «Descripción de la
% empresa» para el caso de la Pasantía Larga/Intermedia. Mínimo dos páginas
% por capítulo (norma 2.6). Redactado en tercera persona.
%
% Los datos de identificación, actividad y modelo de trabajo de la empresa son
% información pública/institucional de Domotai C.A., compartida con otro
% informe de pasantía de la misma empresa disponible como referencia de
% formato; la redacción de este capítulo es propia. La sección «Ámbito de la
% pasantía» y el nombre del tutor industrial son específicos de este trabajo.
%
% PENDIENTE: confirmar con la empresa el nombre y el cargo exactos del tutor
% industrial antes de la entrega final.

\chapter{Descripción de la empresa}

La pasantía se desarrolla en Domotai C.A., una empresa venezolana de desarrollo de software
enfocada en la aplicación de la inteligencia artificial a productos y procesos. Este capítulo
describe su actividad, sus líneas de servicio y su forma de trabajo, y sitúa el proyecto de la
pasantía dentro de ese contexto.

\section{Identificación y ubicación}

Domotai C.A. opera desde el Edificio Caracas Campus, en la calle Altagracia, urbanización La
Trinidad, municipio Baruta del estado Miranda, en la ciudad de Caracas. Su actividad no se limita
al mercado local: atiende clientes de sectores diversos bajo un modelo de trabajo que no exige la
presencia física del cliente en sus instalaciones.

\section{Actividad y servicios}

La empresa combina el criterio de sus profesionales con herramientas de inteligencia artificial
para el desarrollo de productos de software. Su oferta se organiza en las líneas recogidas en la
tabla \ref{tab:servicios-domotai}.

\begin{table}[H]
  \centering
  \caption{Líneas de servicio de Domotai C.A.}
  \label{tab:servicios-domotai}
  \begin{tabular}{@{}p{0.32\textwidth}p{0.56\textwidth}@{}}
    \toprule
    \textbf{Línea} & \textbf{Alcance} \\
    \midrule
    Desarrollo web & Sitios corporativos y aplicaciones de cliente desarrollados a medida \\
    Desarrollo móvil & Aplicaciones multiplataforma \\
    Automatización con inteligencia artificial & Agentes y flujos automatizados integrados en
      los procesos del cliente \\
    Integración \emph{hardware}-\emph{software} & Soluciones que articulan componentes físicos
      con sistemas informáticos \\
    Diseño de interfaz y experiencia de usuario & Definición de la interacción y del aspecto
      visual de los productos \\
    Asesoría tecnológica & Acompañamiento en decisiones de arquitectura y selección de
      herramientas \\
    \bottomrule
  \end{tabular}
\end{table}

De estas líneas, la automatización mediante inteligencia artificial es la que enmarca
directamente el proyecto de esta pasantía y a la que se vuelve al cierre del capítulo.

\section{Modelo de trabajo}

Domotai C.A. acumula más de veinte proyectos completados, con un plazo de entrega característico
del orden de seis semanas: un modo de operar orientado a ciclos cortos, en los que el alcance del
producto se acota para poder entregarse dentro de ese plazo en lugar de plantearse como un
desarrollo prolongado. Su cartera abarca dominios heterogéneos ---diseño y arte, hostelería,
interiorismo, comercio electrónico---, lo que indica que su especialización no reside en un sector
vertical concreto, sino en un conjunto de capacidades técnicas aplicables a dominios distintos.

La contratación se articula bajo dos modalidades: por proyecto, con un alcance cerrado de
antemano, y por suscripción, que da al cliente acceso continuado a la capacidad de desarrollo de
la empresa. La segunda modalidad supone una relación sostenida en el tiempo, lo que favorece la
acumulación de conocimiento sobre el dominio de cada cliente.

\section{Ámbito de la pasantía}

El proyecto desarrollado durante esta pasantía se sitúa en la línea de automatización con
inteligencia artificial descrita en la sección \ref{tab:servicios-domotai}, aplicada esta vez al
dominio de la ciberseguridad: un motor que triaje alertas de seguridad de red, las clasifique y
priorice, y justifique cada decisión mediante un modelo de lenguaje con generación aumentada por
recuperación, en lugar de mediante trabajo manual especializado de un analista.

El encaje con la empresa es directo por dos razones. La primera es de mercado: Domotai C.A. no
compite en un segmento vertical fijo, sino que ofrece capacidades de automatización aplicables a
dominios distintos, y la seguridad de red ---con su necesidad de reducir el trabajo manual repetitivo
de un especialista--- es uno de esos dominios. La segunda es de restricción compartida: al igual que
en el resto de su cartera, la solución debía operar sin depender de servicios en la nube de
terceros, de forma que los datos del cliente ---en este caso, sus alertas de seguridad--- no salgan de
su propia red.

La supervisión de la pasantía corresponde a David Enrique Altuve Villalba, Gerente de Operaciones
de la compañía, designado tutor industrial. La estancia se desarrolla en modalidad presencial en la
sede de la empresa, con la duración de veinte semanas que corresponde a una pasantía larga de la
Universidad Simón Bolívar.
